% Generated draft: requires mathematical and source review.
\documentclass[11pt]{article}
\usepackage[margin=25mm]{geometry}
\usepackage{fontspec}
\setmainfont{Latin Modern Roman}
\usepackage{amsmath,amssymb,mathtools,unicode-math}
\setmathfont{Latin Modern Math}
\usepackage{enumitem,needspace,xurl,hyperref}
\setlength{\parindent}{0pt}
\setlength{\parskip}{5pt}
\setlength{\emergencystretch}{3em}
\newcommand{\N}{\mathbb{N}}
\newcommand{\Z}{\mathbb{Z}}
\newcommand{\Q}{\mathbb{Q}}
\newcommand{\R}{\mathbb{R}}
\newcommand{\C}{\mathbb{C}}
\newcommand{\F}{\mathbb{F}}
\newcommand{\A}{\mathbb{A}}
\newcommand{\PP}{\mathbb{P}}
\newcommand{\E}{\mathbb{E}}
\newcommand{\eps}{\varepsilon}
\newcommand{\dd}{\,\mathrm{d}}
\begin{document}
\section*{UnsolvedMath ID: 30003061}
\textbf{Source problem code:} catalogue entry
\subsection{Definitions and mathematical statement}
The supplied question omits hypotheses on its algebraic data. The following gives a conventional framework for interpreting its notation; these hypotheses are not independently confirmed by the supplied record.

Let \(k\) be a field of characteristic \(p>0\), let \(G\) be a finite group with \(p\mid |G|\), and let \(A\) be an associative unital \(k\)-algebra equipped with an action of \(G\) by \(k\)-algebra automorphisms. Define the smash product \(B=A\# kG\) to have underlying vector space \(A\otimes_k kG\), with multiplication
\[
(a\# g)(b\# h)=a\,g(b)\# gh
\qquad(a,b\in A,\ g,h\in G).
\]
Suppose that \(M\) is a \(B\)-bimodule, and regard it as an \(A\)-bimodule by restriction along \(a\mapsto a\#1\). For an algebra \(C\) and a \(C\)-bimodule \(N\), write
\[
C^e=C\otimes_k C^{\mathrm{op}},\qquad
\operatorname{HH}^n(C,N)=\operatorname{Ext}_{C^e}^n(C,N),\qquad
\operatorname{HH}^*(C,N)=\bigoplus_{n\geq0}\operatorname{HH}^n(C,N).
\]
Writing \(u_g=1_A\#g\), the action \(g\cdot m=u_g m u_{g^{-1}}\), together with the action on \(A\), induces the usual action on Hochschild cohomology. In cochain degree \(n\), this action is represented by
\[
(g\cdot f)(a_1,\ldots,a_n)
=u_g f\bigl(g^{-1}(a_1),\ldots,g^{-1}(a_n)\bigr)u_{g^{-1}}.
\]
Define
\[
\operatorname{HH}^*(A,M)^G
=\{x\in\operatorname{HH}^*(A,M):g\cdot x=x\text{ for every }g\in G\}.
\]

\textbf{Source question, with unresolved qualifications.} In the modular case \(p\mid |G|\), determine the relationship between
\[
\operatorname{HH}^*(A,M)^G
\quad\text{and}\quad
\operatorname{HH}^*(A\# kG,M).
\]
The source calls these objects graded algebras. For an arbitrary bimodule \(M\), however, the definitions above supply graded vector spaces without a specified internal cup product. Suitable compatible multiplicative structure on \(M\), or an intended specialization of \(M\), must be supplied before the graded-algebra formulation is determined. The source does not specify this structure, the full class of admissible \(k,G,A,M\), or a particular comparison map, identity, inequality, or other relationship to be established. Accordingly, this is an unresolved comparison question in a conditional framework, not a fully specified conjectural assertion.
\subsection{Short English statement}
For a finite group acting on an algebra over a field whose characteristic divides the group order, compare the invariant Hochschild cohomology of the algebra with the Hochschild cohomology of its smash product. The source leaves the coefficient object \(M\), the structure needed to obtain graded algebras, and the precise intended comparison insufficiently specified.
\subsection{Sources}
\begin{itemize}
\item Supplied local UnsolvedMath registry record, permanent ID \(30003061\), title \emph{Invariant and Smash-Product Hochschild Cohomology}; exact source problem code \texttt{catalogue entry}. The record places the problem in the Oberwolfach Reports Open Problems collection and gives the proposed year as 2016; no individual proposer is supplied.
\item Supplied report citation: \emph{Hochschild Cohomology in Algebra, Geometry, and Topology} (2016), Oberwolfach Reports, pp.\ 449--506, DOI \texttt{10.4171/owr/2016/10}. \url{https://doi.org/10.4171/owr/2016/10}. The background describes Oberwolfach Reports as documenting mathematical workshops and their open-problem discussions. No problem-specific page or named contributor is supplied.
\item Supplied literature source, labelled ``OWR report'': \url{https://doi.org/10.4171/owr/2016/10}.
\item Supplied literature source, labelled ``smash products in the nonsemisimple case'': \url{https://arxiv.org/abs/1511.05388}.
\item Supplied literature source, labelled ``smash-product spectral sequences'': \url{https://arxiv.org/abs/1401.3551}.
\item Literature assessment dated exactly \texttt{2026-08-21}: the supplied curation labels the problem partially solved and states that a modular spectral sequence gives inequalities and collapse criteria for cyclic \(p\)-groups and certain extensions, while the general relation between invariant Hochschild cohomology and smash-product Hochschild cohomology remains unresolved. This is supplied dated curation, not an independently verified literature assessment. The background expressly requests independent verification of time-sensitive status and cited claims.
\item The registry labels the statement \texttt{corrected\_verified} and says that mathematical content was preserved while notation, typography, or omitted definitions were repaired. This is supplied registry metadata, not independent mathematical verification of this draft.
\end{itemize}
\end{document}
